% Replacement text prepared from the implementation at commit 6a82c276.
% The bibliography keys Krumholz2007 and Zhang2011 should be replaced by the
% corresponding keys in the main manuscript bibliography if they differ.

\section{Governing Equations and FLD Closure}
\label{sec:basic_equations}

\subsection{Mixed-frame grey radiation hydrodynamics}
\label{sec:rhd_equations}

We solve the non-relativistic, frequency-integrated grey
radiation-hydrodynamics equations without magnetic fields or gravity.  The
material variables are the mass density $\rho$, velocity $\bm v$, gas pressure
$P_{\rm gas}$, and total gas energy density
\begin{align}
    e_{\rm gas,tot}
    = e_{\rm gas}+\frac{1}{2}\rho |\bm v|^2,
    \label{eq:gas_total_energy_definition}
\end{align}
where $e_{\rm gas}$ is the gas internal energy density.  The radiation field
is represented by the lab-frame radiation energy density $E_{\rm rad}$.  The
opacities, on the other hand, are evaluated in the frame comoving with the
gas.  Thus, the formulation is a mixed-frame FLD system
\citep{Krumholz2007,Zhang2011}.

The implementation evolves the following reduced system:
\begin{align}
    \frac{\partial \rho}{\partial t}
    +\nabla\cdot(\rho\bm v)
    &=0,
    \label{eq:rhd_mass}\\
    \frac{\partial(\rho\bm v)}{\partial t}
    +\nabla\cdot
    \left(\rho\bm v\bm v+P_{\rm gas}{\sf I}\right)
    &=-\lambda\nabla E_{\rm rad},
    \label{eq:rhd_momentum}\\
    \frac{\partial e_{\rm gas,tot}}{\partial t}
    +\nabla\cdot
    \left[(e_{\rm gas,tot}+P_{\rm gas})\bm v\right]
    &=-Q_{\rm th}+W_{\rm mf},
    \label{eq:rhd_gas_energy}\\
    \frac{\partial E_{\rm rad}}{\partial t}
    +\nabla\cdot\left(\zeta E_{\rm rad}\bm v\right)
    &=\nabla\cdot\left(D\nabla E_{\rm rad}\right)
      +Q_{\rm th}-W_{\rm mf}.
    \label{eq:rhd_radiation_energy}
\end{align}
Here
\begin{align}
    D&=\frac{c\lambda}{\sigma_{\rm R}},
    \label{eq:fld_diffusion_coefficient}\\
    \zeta&=\frac{3-\chi}{2},
    \label{eq:radiation_advection_factor}\\
    Q_{\rm th}
    &=c\sigma_{\rm P}
      \left(aT_{\rm gas}^{4}-E_{\rm rad}\right),
    \label{eq:thermal_coupling}\\
    W_{\rm mf}
    &=\lambda\left(\frac{2\sigma_{\rm P}}{\sigma_{\rm R}}-1\right)
      \bm v\cdot\nabla E_{\rm rad}.
    \label{eq:mixed_frame_work}
\end{align}
A positive $Q_{\rm th}$ transfers energy from the gas to the radiation field.
The term $W_{\rm mf}$ is the explicit $O(v/c)$ mixed-frame energy correction;
it contains both diffusion work and the Lorentz-transformation correction to
emission and absorption.  It is applied with equal and opposite signs to the
gas and radiation energies.  The lab-frame radiation flux represented by
Equation~\eqref{eq:rhd_radiation_energy} is therefore
\begin{align}
    \bm F_{\rm rad}
    =-D\nabla E_{\rm rad}+\zeta E_{\rm rad}\bm v.
    \label{eq:lab_frame_radiation_flux}
\end{align}
The second term is the advective radiation-enthalpy flux.  In particular,
$\zeta\rightarrow4/3$ in the optically thick limit and
$\zeta\rightarrow1$ in the streaming limit.

The exchange terms cancel exactly upon adding
Equations~\eqref{eq:rhd_gas_energy} and
\eqref{eq:rhd_radiation_energy}.  In the absence of boundary fluxes, the
continuum equations consequently conserve the sum of gas and radiation
energies.  In contrast, FLD does not evolve an independent radiation momentum
density.  The radiation force in Equation~\eqref{eq:rhd_momentum} is therefore
included as a source in the gas momentum equation rather than as an exchange
with an evolved radiation momentum equation.

For an ideal gas,
\begin{align}
    e_{\rm gas}
    &=\frac{\rho k_{\rm B}T_{\rm gas}}
            {(\gamma-1)\mu m_{\rm H}},
    \label{eq:ideal_gas_internal_energy}\\
    P_{\rm gas}&=(\gamma-1)e_{\rm gas}.
    \label{eq:ideal_gas_pressure}
\end{align}
The code also supports a general, tabulated equation of state, in which case
$T_{\rm gas}=T(\rho,e_{\rm gas})$ and
$P_{\rm gas}=P(\rho,e_{\rm gas})$ are supplied by the equation-of-state
module.  The radiation temperature is defined by
\begin{align}
    E_{\rm rad}=aT_{\rm rad}^{4}.
    \label{eq:radiation_temperature}
\end{align}

\subsection{FLD closure}
\label{sec:fld_equations}

We distinguish the mass opacities, $\kappa_{\rm P}$ and $\kappa_{\rm R}$,
from the corresponding interaction coefficients per unit length:
\begin{align}
    \sigma_{\rm P}=\rho\kappa_{\rm P},
    \qquad
    \sigma_{\rm R}=\rho\kappa_{\rm R}.
    \label{eq:opacity_definitions}
\end{align}
The quantities $\kappa_{\rm P}$ and $\kappa_{\rm R}$ have units of
${\rm cm^2\,g^{-1}}$, whereas $\sigma_{\rm P}$ and $\sigma_{\rm R}$ have
units of ${\rm cm^{-1}}$.  The comoving-frame diffusive flux is
\begin{align}
    \bm F_{\rm rad}^{\rm diff}
    =-\frac{c\lambda}{\sigma_{\rm R}}\nabla E_{\rm rad}.
    \label{eq:fld_flux}
\end{align}
We use the Levermore--Pomraning flux limiter
\citep{Levermore1981-wp},
\begin{align}
    \lambda(R)&=\frac{2+R}{6+3R+R^2},
    \label{eq:flux_limiter}\\
    R&=\frac{|\nabla E_{\rm rad}|}
             {\sigma_{\rm R}E_{\rm rad}},
    \label{eq:flux_limiter_argument}\\
    \chi&=\lambda+(\lambda R)^2.
    \label{eq:eddington_factor}
\end{align}
In the optically thick limit, $R\rightarrow0$,
$\lambda\rightarrow1/3$, and $\chi\rightarrow1/3$.  This gives the ordinary
diffusion coefficient $D\rightarrow c/(3\sigma_{\rm R})$.  In the optically
thin limit, $\lambda\rightarrow1/R$ and $\chi\rightarrow1$, so that the
diffusive flux is limited by
$|\bm F_{\rm rad}^{\rm diff}|\rightarrow cE_{\rm rad}$.

The full Levermore closure may be written as
\begin{align}
    {\sf P}_{\rm rad}
    &={\sf f}E_{\rm rad},\\
    {\sf f}
    &=\frac{1-\chi}{2}{\sf I}
      +\frac{3\chi-1}{2}\hat{\bm n}\hat{\bm n},\\
    \hat{\bm n}
    &=\frac{\nabla E_{\rm rad}}{|\nabla E_{\rm rad}|}.
    \label{eq:levermore_eddington_tensor}
\end{align}
However, the present reduced solver neither constructs this tensor nor
evolves a radiation momentum equation.  It uses $\chi$ only in the advection
factor $\zeta$ and uses the scalar FLD pressure $\lambda E_{\rm rad}$ in the
HLLC wave construction.  The corresponding radiation force is evaluated as
\begin{align}
    \frac{\sigma_{\rm R}}{c}\bm F_{\rm rad}^{\rm diff}
    =-\lambda\nabla E_{\rm rad},
    \label{eq:radiation_force}
\end{align}
which is the source appearing in Equation~\eqref{eq:rhd_momentum}.

For tests in which a fixed limiter is requested, the implementation sets both
$\lambda$ and $\chi$ to $1/3$ everywhere.  The radiation-force coupling and
the mixed-frame correction can also be disabled independently.  Unless
otherwise stated, Equations~\eqref{eq:rhd_mass}--
\eqref{eq:rhd_radiation_energy} refer to the fully coupled choice in which
both terms are retained.

\subsection{Operator splitting used by the implementation}
\label{sec:fld_operator_split}

The update is split into an explicit hyperbolic part and an implicit
diffusion--thermal-coupling part.  In the current driver, one implicit NRFLD
substep is taken at the beginning of each global timestep, followed by the
stages of the explicit hydrodynamic integrator.  The explicit subsystem is
\begin{align}
    \frac{\partial \rho}{\partial t}
    +\nabla\cdot(\rho\bm v)&=0,\\
    \frac{\partial(\rho\bm v)}{\partial t}
    +\nabla\cdot
      \left(\rho\bm v\bm v+P_{\rm gas}{\sf I}\right)
      &=-\lambda\nabla E_{\rm rad},\\
    \frac{\partial e_{\rm gas,tot}}{\partial t}
    +\nabla\cdot
      \left[(e_{\rm gas,tot}+P_{\rm gas})\bm v\right]
      &=W_{\rm mf},\\
    \frac{\partial E_{\rm rad}}{\partial t}
    +\nabla\cdot(\zeta E_{\rm rad}\bm v)
      &=-W_{\rm mf}.
    \label{eq:explicit_subsystem}
\end{align}
Gas and radiation face states are reconstructed at the same spatial order.
The HLLC-FLD Riemann solver uses
\begin{align}
    P_{\rm tot}=P_{\rm gas}+\lambda E_{\rm rad}
    \label{eq:hllc_total_pressure}
\end{align}
to construct the common Godunov state and estimates the acoustic speed as
\begin{align}
    c_{\rm eff}^{2}
    =c_{\rm gas}^{2}
     +\frac{\lambda(E_{\rm rad}+\lambda E_{\rm rad})}{\rho}.
    \label{eq:radiation_modified_sound_speed}
\end{align}
This wave-speed estimate corresponds to the approximate characteristic
relation $\zeta\simeq1+\lambda$, which is exact in both the diffusion and
streaming limits and remains close for the adopted limiter.  Internally, the
Riemann solver first constructs an HLLC total-energy flux using
$e_{\rm gas,tot}+E_{\rm rad}$ and $P_{\rm tot}$.  It assigns
$\zeta E_{\rm rad}\bm v$ to the radiation energy and returns the remainder as
the gas-energy flux, so that the gas and radiation updates share one
conservative interface flux.

The scalar radiation-pressure contribution is subsequently removed from the
returned gas momentum flux, and the equivalent force
$-\lambda\nabla E_{\rm rad}$ is applied explicitly using the Godunov
radiation energies on the cell faces.  The radiation advection flux uses the
same HLLC contact speed and is
$\zeta E_{\rm rad}\bm v$.

During the implicit substep, $\rho$ and $\bm v$ are fixed and only
$e_{\rm gas}$ and $E_{\rm rad}$ are updated:
\begin{align}
    \frac{\partial e_{\rm gas}}{\partial t}
    &=-Q_{\rm th},
    \label{eq:implicit_gas_energy}\\
    \frac{\partial E_{\rm rad}}{\partial t}
    &=\nabla\cdot(D\nabla E_{\rm rad})+Q_{\rm th}.
    \label{eq:implicit_radiation_energy}
\end{align}
Thus, neither ${\sf P}_{\rm rad}:\nabla\bm v$ nor $W_{\rm mf}$ is part of the
current Newton--multigrid solve.  The former is represented by the explicit
radiation-enthalpy transport and radiation-force coupling, whereas the latter
is applied directly as an explicit, equal-and-opposite gas--radiation energy
exchange.


\section{Method}
\label{sec:method}

We implemented the above FLD system in the open-source astrophysical
hydrodynamics code Athena++ \citep{Stone2020-np}.  The current NRFLD
configuration is Newtonian and non-magnetic, and uses the dedicated HLLC-FLD
Riemann solver.  The Newton--multigrid diffusion solve currently requires a
three-dimensional uniform mesh and logically cubic MeshBlocks; static mesh
refinement and adaptive mesh refinement are supported through the multilevel
boundary and multigrid operators.

\subsection{Implicit FLD update}
\label{sec:method_overview}

Let the superscript $-$ denote the state entering an implicit substep and the
superscript $+$ the state leaving it.  A backward-Euler update solves for
$e_{\rm gas}^{+}$ and $E_{\rm rad}^{+}$.  The nonlinear residuals implemented
in the code are
\begin{align}
    F_{\rm gas}
    &=e_{\rm gas}^{+}-e_{\rm gas}^{-}
      +\Delta t\,c\sigma_{\rm P}
       \left[aT^4(e_{\rm gas}^{+},\rho)-E_{\rm rad}^{+}\right],
    \label{eq:Fgas_residual}\\
    F_{\rm rad}
    &=E_{\rm rad}^{+}-E_{\rm rad}^{-}
      -\Delta t\Bigg\{
       c\sigma_{\rm P}
       \left[aT^4(e_{\rm gas}^{+},\rho)-E_{\rm rad}^{+}\right]
       +\nabla\cdot\left(D^{-}\nabla E_{\rm rad}^{+}\right)
       \Bigg\}.
    \label{eq:Frad_residual}
\end{align}
The Planck and Rosseland interaction coefficients and the face-centered
diffusion coefficient $D^{-}$ are frozen at the beginning of the implicit
substep.  In particular, $R$, $\lambda$, and the face diffusion coefficients
are evaluated from the pre-update radiation field and are not recomputed
during the Newton iterations.  The thermal source remains nonlinear through
$T(e_{\rm gas}^{+},\rho)^4$.

For a Newton correction $(\delta e_{\rm gas},\delta E_{\rm rad})$, the local
gas-energy correction is eliminated analytically:
\begin{align}
    \delta e_{\rm gas}
    =-\frac{F_{\rm gas}
       +(\partial F_{\rm gas}/\partial E_{\rm rad})\delta E_{\rm rad}}
      {\partial F_{\rm gas}/\partial e_{\rm gas}}.
    \label{eq:local_gas_elimination}
\end{align}
The resulting scalar elliptic equation for $\delta E_{\rm rad}$ is solved by
the linear multigrid solver.  Its default smoother is red--black weighted
Jacobi (the input value \texttt{jacobi-rb}), with trilinear prolongation.
After each linear solve, the Newton update is subject to positivity floors and
a global backtracking line search when needed.  The converged radiation and
gas internal energies are then copied back to the radiation and hydrodynamic
state variables.
